% ======================================================================
%  AT-MONO110 — Chapter 11: The Standard Model
%  Canonical Monograph (v2.0) · The Actualization Theory:
%  A Reconstruction of Physics from Difference, Actualization and Spectrum
%
%  Status: COMPLETE — PASS-READY (MONO007 referee-ready architecture)
%  Inputs: Chapter01_Difference.tex ... Chapter10_GravityAndSpacetime.tex
%  Method: assembly only — canonical results only, no new physics, no speculation
%  Citations: QG149 (sector exponents), QG150 (mode access), QG157 (effective access),
%             QG161 (gauge origin), QG162 (coupling origin), QG165 (CKM),
%             QG167 (PMNS), QG168 (weak boson masses), QG169 (Higgs),
%             QG172 (neutrino masses), QG176 (Higgs blind reconstruction),
%             QG178 (electron g-2), QG179 (Majorana), QG180 (oblique)
%
%  Usage: % ======================================================================
%  AT-MONO110 — Chapter 11: The Standard Model
%  Canonical Monograph (v2.0) · The Actualization Theory:
%  A Reconstruction of Physics from Difference, Actualization and Spectrum
%
%  Status: COMPLETE — PASS-READY (MONO007 referee-ready architecture)
%  Inputs: Chapter01_Difference.tex ... Chapter10_GravityAndSpacetime.tex
%  Method: assembly only — canonical results only, no new physics, no speculation
%  Citations: QG149 (sector exponents), QG150 (mode access), QG157 (effective access),
%             QG161 (gauge origin), QG162 (coupling origin), QG165 (CKM),
%             QG167 (PMNS), QG168 (weak boson masses), QG169 (Higgs),
%             QG172 (neutrino masses), QG176 (Higgs blind reconstruction),
%             QG178 (electron g-2), QG179 (Majorana), QG180 (oblique)
%
%  Usage: % ======================================================================
%  AT-MONO110 — Chapter 11: The Standard Model
%  Canonical Monograph (v2.0) · The Actualization Theory:
%  A Reconstruction of Physics from Difference, Actualization and Spectrum
%
%  Status: COMPLETE — PASS-READY (MONO007 referee-ready architecture)
%  Inputs: Chapter01_Difference.tex ... Chapter10_GravityAndSpacetime.tex
%  Method: assembly only — canonical results only, no new physics, no speculation
%  Citations: QG149 (sector exponents), QG150 (mode access), QG157 (effective access),
%             QG161 (gauge origin), QG162 (coupling origin), QG165 (CKM),
%             QG167 (PMNS), QG168 (weak boson masses), QG169 (Higgs),
%             QG172 (neutrino masses), QG176 (Higgs blind reconstruction),
%             QG178 (electron g-2), QG179 (Majorana), QG180 (oblique)
%
%  Usage: % ======================================================================
%  AT-MONO110 — Chapter 11: The Standard Model
%  Canonical Monograph (v2.0) · The Actualization Theory:
%  A Reconstruction of Physics from Difference, Actualization and Spectrum
%
%  Status: COMPLETE — PASS-READY (MONO007 referee-ready architecture)
%  Inputs: Chapter01_Difference.tex ... Chapter10_GravityAndSpacetime.tex
%  Method: assembly only — canonical results only, no new physics, no speculation
%  Citations: QG149 (sector exponents), QG150 (mode access), QG157 (effective access),
%             QG161 (gauge origin), QG162 (coupling origin), QG165 (CKM),
%             QG167 (PMNS), QG168 (weak boson masses), QG169 (Higgs),
%             QG172 (neutrino masses), QG176 (Higgs blind reconstruction),
%             QG178 (electron g-2), QG179 (Majorana), QG180 (oblique)
%
%  Usage: \input{Chapter11_StandardModel.tex} after the preamble
%         that defines the theorem environments (or include via the master file).
% ======================================================================

% ---- Theorem environments (define in the master preamble if not present) ----
% \newtheorem{definition}{Definition}[chapter]
% \newtheorem{lemma}[definition]{Lemma}
% \newtheorem{theorem}[definition]{Theorem}
% \newtheorem{corollary}[definition]{Corollary}
% \newtheorem{remark}[definition]{Remark}

\chapter{The Standard Model}
\label{c11:chap:sm}

\section{Introduction}
\label{c11:sec:sm-introduction}

The preceding chapters established the foundation, the spectrum, the operator basis, and
the quantum and gravitational structure of The Actualization Theory. This chapter
reconstructs the \emph{standard-model observables} from the same foundation --- the fermion
sectors, the gauge structure, the couplings, the flavor mixing, the weak scale and the
Higgs, and a set of precision predictions --- as derived reads of the D96 spectral moments
\cite{c11:QG149,c11:QG150,c11:QG157}.

The central results are: the fermion sectors are reads of the spectral access counts
\cite{c11:QG149,c11:QG150}; the gauge structure $1+3+8$ is the degree-12 structure of the D96
sector \cite{c11:QG161}; the couplings are ratios of spectral constants \cite{c11:QG162}; the
flavor mixing matrices are octave-transition reads \cite{c11:QG165,c11:QG167}; the weak scale and
the Higgs mass are spectral constructions \cite{c11:QG168,c11:QG169}; the neutrino masses are
closed-form D96 laws \cite{c11:QG172}; and a set of precision predictions --- the oblique
parameters $S,T,U$, the electron g-2, and the Majorana character --- follow
\cite{c11:QG176,c11:QG178,c11:QG179,c11:QG180}.

Throughout, the hosted-standard-model caveat is retained: every \emph{observable} is
derived from the D96 moments, while the standard-model \emph{Lagrangian} and its dynamics
are hosted --- a boundary, not a derivation gap, per the referee-safe architecture of
MONO007 \cite{c11:MONO007}.

The dimensionless D96 laws are then calibrated to measured anchors: the weak scale $v$
for the weak/Higgs sector, the electron mass $m_e$ for the absolute quark masses, and the
standard unit conversion when a result is reported in SI or GeV units. That calibration
step fixes the physical units; it is not a new dynamical input.

\subsection{Derived Observables versus Hosted Dynamics}
\label{c11:sec:derived-vs-hosted}

\begin{quote}
\textbf{Derived:} primitives $\rightarrow$ actualization $\rightarrow$ spectrum $\rightarrow$
observables.\\
\textbf{Hosted:} SM Lagrangian, interaction dynamics, and external unit conventions.
\end{quote}

What is reconstructed here is the observable content of the standard model: masses,
couplings, mixings, and precision numbers. What is not reconstructed is the full dynamical
Lagrangian/vertex structure; that remains hosted and is treated as a boundary condition of
the monograph, not as a derived output.

\section{Fermion Sector}
\label{c11:sec:fermion-sector}

\begin{theorem}[The fermion sectors are spectral access reads]
\label{c11:thm:fermion-access}
The \emph{fermion sector} of the theory --- the set of fermion families and hierarchies ---
is a spectral access read: each sector couples to the spectrum through a distinct moment of
the multiplicity distribution, the neutral, full, doublet-occupancy, and octave-occupation
accesses being the half, first, second, and octave moments.
\end{theorem}

\begin{proof}
The mode-access origin \cite{c11:QG150} and the effective-access program \cite{c11:QG157}
establish that each sector couples through a distinct moment of the multiplicity
distribution: the neutral access is the half-moment $\Sigma\sqrt{m} = 64.08$, the full
access the first moment $\Sigma m = 95$, the doublet-occupancy access the second
moment $\Sigma m^2 = 229$, and the octave-occupation access
$\mathrm{occMom} = 1900.25$. The sector exponents follow from these access counts
\cite{c11:QG149}, and the mass ratios are functions of the exponents; since the counts are
themselves derived from the spectrum (Chapter~\ref{c06:chap:d96},
Corollary~\ref{c06:cor:moment-ladder}, Corollary~\ref{c06:cor:constants-derived}), the
fermion sectors and their mass hierarchy are spectral access reads.
\end{proof}

\section{Gauge Structure}
\label{c11:sec:gauge}

\begin{theorem}[The gauge structure is the D96 degree-12 observable sector]
\label{c11:thm:gauge-structure}
The gauge structure of the standard model, $U(1)\times SU(2)\times SU(3)$ --- the
$1+3+8$ gauge sector --- is reconstructed from the D96 sector structure: the degree-12 content of
the $C_{96}$ generator ring.
\end{theorem}

\begin{proof}
The gauge-origin audit \cite{c11:QG161} establishes that the $1+3+8$ gauge sector is the
degree-12 structure of the $C_{96}$ sector: the gauge generators are the D96 spectral
content, giving one $U(1)$, three $SU(2)$, and eight $SU(3)$ generators --- reconstructed as
an observable sector from the D96 spectrum, not imported, and introducing no new primitive.
Since the D96 spectrum is itself the derived output of the attractor
(Chapter~\ref{c06:chap:d96}, Theorem~\ref{c06:thm:d96-derived}), the gauge structure is
twice-emergent as an observable sector; the interaction dynamics remain hosted.
\end{proof}

\section{Coupling Constants}
\label{c11:sec:couplings}

\begin{theorem}[The couplings are spectral ratios]
\label{c11:thm:couplings}
The gauge couplings are derived as ratios of D96 spectral constants:
\begin{equation}
\label{c11:eq:couplings}
\frac{1}{\alpha_{\mathrm{em}}} = \Sigma m + \#d = 137,
\qquad
\alpha_{\mathrm{weak}} = \frac{3}{\Sigma m},
\qquad
\alpha_{\mathrm{strong}} = \frac{8}{\Sigma\sqrt{m}},
\qquad
\sin^2\theta_W = 0.2316.
\end{equation}
\end{theorem}

\begin{proof}
The coupling-origin audit \cite{c11:QG162} establishes the constructions: the fine-structure
inverse is the total mode count plus the doublet count, $1/\alpha_{\mathrm{em}} = 137$; the
weak and strong couplings are ratios of the spectral constants; the Weinberg angle is a
derived ratio. Each coupling is therefore a ratio of spectral constants ($\Sigma m$,
$\Sigma\sqrt{m}$, $\#d$) --- derived, not fitted: the constants are themselves derived
from the spectrum (Chapter~\ref{c06:chap:d96}, Corollary~\ref{c06:cor:constants-derived}),
and no fitted parameters or imported values enter.
\end{proof}

\section{Flavor Mixing}
\label{c11:sec:mixing}

\begin{theorem}[The CKM matrix is an octave-transition read]
\label{c11:thm:ckm}
The CKM matrix is derived as the quark mixing read of the spectral structure: the elements
are ratios of doublet, octave, and occupancy quantities.
\end{theorem}

\begin{proof}
The CKM-origin audit \cite{c11:QG165} establishes the constructions:
$V_{us} = \#d/(2\Sigma m)$, $V_{cb} = (\omega_0/\omega_2)^{\delta_d}$,
$V_{ub} = 2V_{cb}\cdot(\mathrm{occ}_0/\mathrm{occ}_2)$, with mean deviation $0.58\%$; the
CKM CP phase and the Jarlskog invariant follow from the chiral circulation \cite{c11:QG166}.
\end{proof}

\begin{theorem}[The PMNS matrix is a T3-only read]
\label{c11:thm:pmns}
The PMNS matrix is derived as the neutrino mixing read: the angles
$\theta_{12} = 33.35^\circ$, $\theta_{23} = 49.72^\circ$, $\theta_{13} = 8.34^\circ$, and
$\delta_\nu = 66.4^\circ$ are the T3-only access reads, with mean deviation $1.5\%$.
\end{theorem}

\begin{proof}
The PMNS-origin audit \cite{c11:QG167} establishes the angles and the CP phase as the
T3-only access reads of the D96 content; hence the PMNS matrix is a derived read, not an
imported input.
\end{proof}

Both the quark (CKM) and the lepton (PMNS) mixing matrices are therefore derived spectral
reads, with no imported mixing input.

\section{Higgs and Weak Scale}
\label{c11:sec:higgs}

\begin{theorem}[The weak scale and the Higgs mass]
\label{c11:thm:weak-higgs}
The weak scale and the Higgs mass are reconstructed from the spectral structure:
\begin{equation}
\label{c11:eq:higgs}
v = (\Sigma m+\#d)\ln(\mathrm{span}) = 254\ \mathrm{GeV},
\qquad
M_H = \sigma_{\mathrm{occ}}\cdot\frac{\mathrm{span}}{2} = 125.25\ \mathrm{GeV},
\end{equation}
with the $W$ and $Z$ masses following from the Weinberg projection.
\end{theorem}

\begin{proof}
The weak-boson-mass origin \cite{c11:QG168} establishes the weak scale
$v = (\Sigma m+\#d)\ln(\mathrm{span}) = 254$ GeV and the $W$ and $Z$ masses
($M_W = 80.1$, $M_Z = 91.4$ GeV) with $\rho = 1$ exactly; the Higgs-mass origin
\cite{c11:QG169} gives $M_H = \sigma_{\mathrm{occ}}\cdot(\mathrm{span}/2) = 125.25$ GeV,
matching observation within $0.003\%$; and the blind reconstruction \cite{c11:QG176}
confirms this mass from the pre-Higgs D96 content only, with no Higgs observable as input.
The D96 factors are dimensionless; the GeV unit is carried by the calibrated electroweak
scale $v$.
\end{proof}

The Higgs is therefore a collective scalar observable of the theory: its mass is a spectral
construction $M_H = \sigma_{\mathrm{occ}}\cdot(\mathrm{span}/2)$, whose constants are
derived from the spectrum (Chapter~\ref{c06:chap:d96},
Corollary~\ref{c06:cor:constants-derived}); it introduces no new primitive \cite{c11:QG169}.
Read strictly, the D96 factor is dimensionless and the GeV unit is inherited from the
calibrated electroweak mass scale. The Higgs sector observable is reconstructed; the Yukawa
and full interaction dynamics remain hosted.

\begin{lemma}[The neutrino masses are closed-form D96 laws]
\label{c11:lem:neutrino-masses}
The neutrino mass splittings are derived as closed-form D96 laws:
$\Delta m^2_{21} = (1/\Sigma\sqrt{m})^2/(\mathrm{span}/2)$ and
$\Delta m^2_{31} = \sin^2\theta_W/\Sigma m$.
\end{lemma}

\begin{proof}
The neutrino-mass origin \cite{c11:QG172} establishes the closed-form laws: the solar and
atmospheric splittings are ratios of the D96 spectral constants, giving
$\Delta m^2_{21} = 7.607\times10^{-5}$ eV$^2$ and
$\Delta m^2_{31} = 2.44\times10^{-3}$ eV$^2$, with $m_2 = 8.72$ meV,
$m_3 = 49.4$ meV, and $\Sigma m_\nu = 0.058$ eV. Hence the neutrino masses are closed-form
D96 laws, not oscillation fits.
\end{proof}

\section{Precision Predictions}
\label{c11:sec:precision}

\begin{theorem}[The oblique parameters]
\label{c11:thm:oblique}
The oblique parameters are derived from the spectral structure:
\begin{equation}
\label{c11:eq:oblique}
S = \frac{\mathrm{occ}_0}{\Sigma m} = \frac{4}{95} = 0.0421,
\qquad
T = 2S = \frac{8}{95} = 0.0842,
\qquad
U = 0,
\end{equation}
with $T = 2S$ exactly and $U = 0$ via the exact tree-level $\rho = 1$.
\end{theorem}

\begin{proof}
The oblique-origin audit \cite{c11:QG180} establishes the constructions: $S$ is the lightest-
octave fraction of the spectrum, $T = 2S$ exactly, and $U = 0$ via the exact SM tree-level
$\rho = 1$, matching the EW global fit within $5.3\%$.
\end{proof}

\begin{theorem}[The electron g-2]
\label{c11:thm:electron-g2}
The electron g-2 is derived from the spectral structure:
\begin{equation}
\label{c11:eq:electron-g2}
a_e = \frac{\alpha}{2\pi}\Bigl(1 - \Bigl(\frac{\mathrm{occ}_0}{\Sigma m}\Bigr)^2\Bigr)
= 1.159655\times10^{-3},
\qquad
\Delta a_e = \Bigl(\frac{\alpha}{2\pi}\Bigr)^3 \mathrm{span}^{1/4}
\Bigl(\frac{\mathrm{occ}_0}{\Sigma m}\Bigr)^3 = 1.86\times10^{-13},
\end{equation}
with the correction negative and the anomaly below $10^{-12}$.
\end{theorem}

\begin{proof}
The electron-g-2 origin \cite{c11:QG178} establishes the construction: the electron anomaly is
the Schwinger term modulated by the lightest-octave fraction, with the negative correction
$(1 - (\mathrm{occ}_0/\Sigma m)^2)$ and the anomaly below $10^{-12}$; the same D96
mechanism produces the muon g-2 \cite{c11:QG171}. The electron anomaly matches experiment
within $0.0003\%$.
\end{proof}

\begin{theorem}[The Majorana character and \texorpdfstring{$0\nu\beta\beta$}{0nubetabeta}]
\label{c11:thm:majorana}
The neutrino is Majorana: it is self-conjugate in its $T_3$-only channel, has no conserved
charge, and has a real mass matrix. The effective Majorana mass is
\begin{equation}
\label{c11:eq:mbb}
m_{\beta\beta} = \bigl|m_1c_{12}^2c_{13}^2 + m_2s_{12}^2c_{13}^2e^{i\alpha_2}
+ m_3s_{13}^2e^{-2i\delta_\nu}\bigr| = 2.02\times10^{-3}\ \mathrm{eV},
\end{equation}
within experimental limits and in reach of next-generation experiments.
\end{theorem}

\begin{proof}
The Majorana-origin audit \cite{c11:QG179} establishes the neutrino character: the $T_3$-only
channel is self-conjugate ($48/95$ modes), the unique $Q=0$ sector provides no conserved
charge, and the reflection automorphism gives a real mass matrix. From the D96 PMNS angles
and masses, the effective Majorana mass is $m_{\beta\beta} = 2.02\times10^{-3}$ eV, within
limits and in reach of next-generation experiments.
\end{proof}

The precision predictions --- $S,T,U$, the electron g-2, and the Majorana character ---
are therefore all derived from the D96 spectral structure; their experimental validation
remains open --- the current frontier, not a derivation gap (VALID001).

\begin{remark}[Hosted-SM caveat]
\label{c11:remark:hosted-sm}
Consistent with MONO007 \cite{c11:MONO007}, the standard-model \emph{Lagrangian} and its
dynamics are hosted --- a boundary, not a derivation claim of this chapter. Every
\emph{observable} presented here is derived from the D96 moments; the Lagrangian form
$L = -\frac{1}{4}F^a_{\mu\nu}F^{a\mu\nu} + i\bar{\psi}\gamma^\mu D_\mu\psi - m\bar{\psi}\psi$
is the actualization-flow action established in Chapter~\ref{c09:chap:qm}
(Lemma~\ref{c09:lem:qm-dynamics}), but the full SM dynamical content remains hosted. No stronger
claim is made.
\end{remark}

\section{Consequences}
\label{c11:sec:sm-consequences}

\begin{theorem}[The standard model is consistent with the canonical architecture]
\label{c11:thm:sm-consistent}
The standard model is consistent with the canonical architecture: every derived SM
observable is a read of the D96 spectral moments, with the hosted dynamics the single
referenced boundary. Dimensional observables are reported relative to calibration anchors,
not from the dimensionless D96 data alone.
\end{theorem}

\begin{proof}
The fermion sector (Theorem~\ref{c11:thm:fermion-access}), the gauge structure
(Theorem~\ref{c11:thm:gauge-structure}), the couplings (Theorem~\ref{c11:thm:couplings}), the flavor
mixing (Theorem~\ref{c11:thm:ckm}, Theorem~\ref{c11:thm:pmns}), the Higgs and weak scale
(Theorem~\ref{c11:thm:weak-higgs}), and the precision predictions
(Theorem~\ref{c11:thm:oblique}, Theorem~\ref{c11:thm:electron-g2}, Theorem~\ref{c11:thm:majorana})
are all derived from the D96 moments. The hosted SM Lagrangian is the single referenced
boundary (Remark~\ref{c11:remark:hosted-sm}).
\end{proof}

\section{Summary}
\label{c11:sec:sm-summary}

This chapter has reconstructed the standard-model observables: the fermion sectors are
spectral access reads (Theorem~\ref{c11:thm:fermion-access}); the $1+3+8$ gauge sector is
the D96 degree-12 observable sector (Theorem~\ref{c11:thm:gauge-structure}); the couplings
are spectral ratios (Theorem~\ref{c11:thm:couplings}); the CKM and PMNS matrices are
derived spectral reads (Theorem~\ref{c11:thm:ckm}, Theorem~\ref{c11:thm:pmns}); and the
weak scale, the Higgs mass, and the neutrino masses are spectral constructions
(Theorem~\ref{c11:thm:weak-higgs}, Lemma~\ref{c11:lem:neutrino-masses}). The precision
predictions --- the oblique parameters $S,T,U$, the electron g-2, and the Majorana
character --- are likewise reconstructed as observable outputs
(Theorem~\ref{c11:thm:oblique}, Theorem~\ref{c11:thm:electron-g2},
Theorem~\ref{c11:thm:majorana}), so the standard model is consistent with the canonical
architecture, with the hosted dynamics the single referenced boundary
(Theorem~\ref{c11:thm:sm-consistent}, Remark~\ref{c11:remark:hosted-sm}). Every SM
observable is therefore a spectral construction, twice-derived from the spectrum
(Chapter~\ref{c06:chap:d96}, Corollary~\ref{c06:cor:constants-derived}); the following
chapter derives cosmology from the same foundation.

\begin{thebibliography}{99}
\bibitem{c11:QG149} AT-QG Phase 149, \emph{Sector Exponents Origin}, Docs/Research.
\bibitem{c11:QG150} AT-QG Phase 150, \emph{Mode Access Origin}, Docs/Research.
\bibitem{c11:QG157} AT-QG Phase 157, \emph{Effective Access Counts}, Docs/Research.
\bibitem{c11:QG161} AT-QG Phase 161, \emph{Gauge Sector Origin}, Docs/Research.
\bibitem{c11:QG162} AT-QG Phase 162, \emph{Gauge Coupling Origin}, Docs/Research.
\bibitem{c11:QG165} AT-QG Phase 165, \emph{CKM Origin}, Docs/Research.
\bibitem{c11:QG166} AT-QG Phase 166, \emph{CKM CP Origin}, Docs/Research.
\bibitem{c11:QG167} AT-QG Phase 167, \emph{PMNS Origin}, Docs/Research.
\bibitem{c11:QG168} AT-QG Phase 168, \emph{Weak Boson Mass Origin}, Docs/Research.
\bibitem{c11:QG169} AT-QG Phase 169, \emph{Higgs Mass Origin}, Docs/Research.
\bibitem{c11:QG171} AT-QG Phase 171, \emph{Muon g-2 Origin}, Docs/Research.
\bibitem{c11:QG172} AT-QG Phase 172, \emph{Neutrino Mass Law}, Docs/Research.
\bibitem{c11:QG176} AT-QG Phase 176, \emph{Higgs Blind Reconstruction}, Docs/Research.
\bibitem{c11:QG178} AT-QG Phase 178, \emph{Electron g-2 Origin}, Docs/Research.
\bibitem{c11:QG179} AT-QG Phase 179, \emph{Majorana Origin}, Docs/Research.
\bibitem{c11:QG180} AT-QG Phase 180, \emph{Oblique Parameters Origin}, Docs/Research.
\bibitem{c11:MONO007} AT-MONO007, \emph{Referee-Readiness Resolution}, Docs/Zenodo/Monograph.
\end{thebibliography}
 after the preamble
%         that defines the theorem environments (or include via the master file).
% ======================================================================

% ---- Theorem environments (define in the master preamble if not present) ----
% \newtheorem{definition}{Definition}[chapter]
% \newtheorem{lemma}[definition]{Lemma}
% \newtheorem{theorem}[definition]{Theorem}
% \newtheorem{corollary}[definition]{Corollary}
% \newtheorem{remark}[definition]{Remark}

\chapter{The Standard Model}
\label{c11:chap:sm}

\section{Introduction}
\label{c11:sec:sm-introduction}

The preceding chapters established the foundation, the spectrum, the operator basis, and
the quantum and gravitational structure of The Actualization Theory. This chapter
reconstructs the \emph{standard-model observables} from the same foundation --- the fermion
sectors, the gauge structure, the couplings, the flavor mixing, the weak scale and the
Higgs, and a set of precision predictions --- as derived reads of the D96 spectral moments
\cite{c11:QG149,c11:QG150,c11:QG157}.

The central results are: the fermion sectors are reads of the spectral access counts
\cite{c11:QG149,c11:QG150}; the gauge structure $1+3+8$ is the degree-12 structure of the D96
sector \cite{c11:QG161}; the couplings are ratios of spectral constants \cite{c11:QG162}; the
flavor mixing matrices are octave-transition reads \cite{c11:QG165,c11:QG167}; the weak scale and
the Higgs mass are spectral constructions \cite{c11:QG168,c11:QG169}; the neutrino masses are
closed-form D96 laws \cite{c11:QG172}; and a set of precision predictions --- the oblique
parameters $S,T,U$, the electron g-2, and the Majorana character --- follow
\cite{c11:QG176,c11:QG178,c11:QG179,c11:QG180}.

Throughout, the hosted-standard-model caveat is retained: every \emph{observable} is
derived from the D96 moments, while the standard-model \emph{Lagrangian} and its dynamics
are hosted --- a boundary, not a derivation gap, per the referee-safe architecture of
MONO007 \cite{c11:MONO007}.

The dimensionless D96 laws are then calibrated to measured anchors: the weak scale $v$
for the weak/Higgs sector, the electron mass $m_e$ for the absolute quark masses, and the
standard unit conversion when a result is reported in SI or GeV units. That calibration
step fixes the physical units; it is not a new dynamical input.

\subsection{Derived Observables versus Hosted Dynamics}
\label{c11:sec:derived-vs-hosted}

\begin{quote}
\textbf{Derived:} primitives $\rightarrow$ actualization $\rightarrow$ spectrum $\rightarrow$
observables.\\
\textbf{Hosted:} SM Lagrangian, interaction dynamics, and external unit conventions.
\end{quote}

What is reconstructed here is the observable content of the standard model: masses,
couplings, mixings, and precision numbers. What is not reconstructed is the full dynamical
Lagrangian/vertex structure; that remains hosted and is treated as a boundary condition of
the monograph, not as a derived output.

\section{Fermion Sector}
\label{c11:sec:fermion-sector}

\begin{theorem}[The fermion sectors are spectral access reads]
\label{c11:thm:fermion-access}
The \emph{fermion sector} of the theory --- the set of fermion families and hierarchies ---
is a spectral access read: each sector couples to the spectrum through a distinct moment of
the multiplicity distribution, the neutral, full, doublet-occupancy, and octave-occupation
accesses being the half, first, second, and octave moments.
\end{theorem}

\begin{proof}
The mode-access origin \cite{c11:QG150} and the effective-access program \cite{c11:QG157}
establish that each sector couples through a distinct moment of the multiplicity
distribution: the neutral access is the half-moment $\Sigma\sqrt{m} = 64.08$, the full
access the first moment $\Sigma m = 95$, the doublet-occupancy access the second
moment $\Sigma m^2 = 229$, and the octave-occupation access
$\mathrm{occMom} = 1900.25$. The sector exponents follow from these access counts
\cite{c11:QG149}, and the mass ratios are functions of the exponents; since the counts are
themselves derived from the spectrum (Chapter~\ref{c06:chap:d96},
Corollary~\ref{c06:cor:moment-ladder}, Corollary~\ref{c06:cor:constants-derived}), the
fermion sectors and their mass hierarchy are spectral access reads.
\end{proof}

\section{Gauge Structure}
\label{c11:sec:gauge}

\begin{theorem}[The gauge structure is the D96 degree-12 observable sector]
\label{c11:thm:gauge-structure}
The gauge structure of the standard model, $U(1)\times SU(2)\times SU(3)$ --- the
$1+3+8$ gauge sector --- is reconstructed from the D96 sector structure: the degree-12 content of
the $C_{96}$ generator ring.
\end{theorem}

\begin{proof}
The gauge-origin audit \cite{c11:QG161} establishes that the $1+3+8$ gauge sector is the
degree-12 structure of the $C_{96}$ sector: the gauge generators are the D96 spectral
content, giving one $U(1)$, three $SU(2)$, and eight $SU(3)$ generators --- reconstructed as
an observable sector from the D96 spectrum, not imported, and introducing no new primitive.
Since the D96 spectrum is itself the derived output of the attractor
(Chapter~\ref{c06:chap:d96}, Theorem~\ref{c06:thm:d96-derived}), the gauge structure is
twice-emergent as an observable sector; the interaction dynamics remain hosted.
\end{proof}

\section{Coupling Constants}
\label{c11:sec:couplings}

\begin{theorem}[The couplings are spectral ratios]
\label{c11:thm:couplings}
The gauge couplings are derived as ratios of D96 spectral constants:
\begin{equation}
\label{c11:eq:couplings}
\frac{1}{\alpha_{\mathrm{em}}} = \Sigma m + \#d = 137,
\qquad
\alpha_{\mathrm{weak}} = \frac{3}{\Sigma m},
\qquad
\alpha_{\mathrm{strong}} = \frac{8}{\Sigma\sqrt{m}},
\qquad
\sin^2\theta_W = 0.2316.
\end{equation}
\end{theorem}

\begin{proof}
The coupling-origin audit \cite{c11:QG162} establishes the constructions: the fine-structure
inverse is the total mode count plus the doublet count, $1/\alpha_{\mathrm{em}} = 137$; the
weak and strong couplings are ratios of the spectral constants; the Weinberg angle is a
derived ratio. Each coupling is therefore a ratio of spectral constants ($\Sigma m$,
$\Sigma\sqrt{m}$, $\#d$) --- derived, not fitted: the constants are themselves derived
from the spectrum (Chapter~\ref{c06:chap:d96}, Corollary~\ref{c06:cor:constants-derived}),
and no fitted parameters or imported values enter.
\end{proof}

\section{Flavor Mixing}
\label{c11:sec:mixing}

\begin{theorem}[The CKM matrix is an octave-transition read]
\label{c11:thm:ckm}
The CKM matrix is derived as the quark mixing read of the spectral structure: the elements
are ratios of doublet, octave, and occupancy quantities.
\end{theorem}

\begin{proof}
The CKM-origin audit \cite{c11:QG165} establishes the constructions:
$V_{us} = \#d/(2\Sigma m)$, $V_{cb} = (\omega_0/\omega_2)^{\delta_d}$,
$V_{ub} = 2V_{cb}\cdot(\mathrm{occ}_0/\mathrm{occ}_2)$, with mean deviation $0.58\%$; the
CKM CP phase and the Jarlskog invariant follow from the chiral circulation \cite{c11:QG166}.
\end{proof}

\begin{theorem}[The PMNS matrix is a T3-only read]
\label{c11:thm:pmns}
The PMNS matrix is derived as the neutrino mixing read: the angles
$\theta_{12} = 33.35^\circ$, $\theta_{23} = 49.72^\circ$, $\theta_{13} = 8.34^\circ$, and
$\delta_\nu = 66.4^\circ$ are the T3-only access reads, with mean deviation $1.5\%$.
\end{theorem}

\begin{proof}
The PMNS-origin audit \cite{c11:QG167} establishes the angles and the CP phase as the
T3-only access reads of the D96 content; hence the PMNS matrix is a derived read, not an
imported input.
\end{proof}

Both the quark (CKM) and the lepton (PMNS) mixing matrices are therefore derived spectral
reads, with no imported mixing input.

\section{Higgs and Weak Scale}
\label{c11:sec:higgs}

\begin{theorem}[The weak scale and the Higgs mass]
\label{c11:thm:weak-higgs}
The weak scale and the Higgs mass are reconstructed from the spectral structure:
\begin{equation}
\label{c11:eq:higgs}
v = (\Sigma m+\#d)\ln(\mathrm{span}) = 254\ \mathrm{GeV},
\qquad
M_H = \sigma_{\mathrm{occ}}\cdot\frac{\mathrm{span}}{2} = 125.25\ \mathrm{GeV},
\end{equation}
with the $W$ and $Z$ masses following from the Weinberg projection.
\end{theorem}

\begin{proof}
The weak-boson-mass origin \cite{c11:QG168} establishes the weak scale
$v = (\Sigma m+\#d)\ln(\mathrm{span}) = 254$ GeV and the $W$ and $Z$ masses
($M_W = 80.1$, $M_Z = 91.4$ GeV) with $\rho = 1$ exactly; the Higgs-mass origin
\cite{c11:QG169} gives $M_H = \sigma_{\mathrm{occ}}\cdot(\mathrm{span}/2) = 125.25$ GeV,
matching observation within $0.003\%$; and the blind reconstruction \cite{c11:QG176}
confirms this mass from the pre-Higgs D96 content only, with no Higgs observable as input.
The D96 factors are dimensionless; the GeV unit is carried by the calibrated electroweak
scale $v$.
\end{proof}

The Higgs is therefore a collective scalar observable of the theory: its mass is a spectral
construction $M_H = \sigma_{\mathrm{occ}}\cdot(\mathrm{span}/2)$, whose constants are
derived from the spectrum (Chapter~\ref{c06:chap:d96},
Corollary~\ref{c06:cor:constants-derived}); it introduces no new primitive \cite{c11:QG169}.
Read strictly, the D96 factor is dimensionless and the GeV unit is inherited from the
calibrated electroweak mass scale. The Higgs sector observable is reconstructed; the Yukawa
and full interaction dynamics remain hosted.

\begin{lemma}[The neutrino masses are closed-form D96 laws]
\label{c11:lem:neutrino-masses}
The neutrino mass splittings are derived as closed-form D96 laws:
$\Delta m^2_{21} = (1/\Sigma\sqrt{m})^2/(\mathrm{span}/2)$ and
$\Delta m^2_{31} = \sin^2\theta_W/\Sigma m$.
\end{lemma}

\begin{proof}
The neutrino-mass origin \cite{c11:QG172} establishes the closed-form laws: the solar and
atmospheric splittings are ratios of the D96 spectral constants, giving
$\Delta m^2_{21} = 7.607\times10^{-5}$ eV$^2$ and
$\Delta m^2_{31} = 2.44\times10^{-3}$ eV$^2$, with $m_2 = 8.72$ meV,
$m_3 = 49.4$ meV, and $\Sigma m_\nu = 0.058$ eV. Hence the neutrino masses are closed-form
D96 laws, not oscillation fits.
\end{proof}

\section{Precision Predictions}
\label{c11:sec:precision}

\begin{theorem}[The oblique parameters]
\label{c11:thm:oblique}
The oblique parameters are derived from the spectral structure:
\begin{equation}
\label{c11:eq:oblique}
S = \frac{\mathrm{occ}_0}{\Sigma m} = \frac{4}{95} = 0.0421,
\qquad
T = 2S = \frac{8}{95} = 0.0842,
\qquad
U = 0,
\end{equation}
with $T = 2S$ exactly and $U = 0$ via the exact tree-level $\rho = 1$.
\end{theorem}

\begin{proof}
The oblique-origin audit \cite{c11:QG180} establishes the constructions: $S$ is the lightest-
octave fraction of the spectrum, $T = 2S$ exactly, and $U = 0$ via the exact SM tree-level
$\rho = 1$, matching the EW global fit within $5.3\%$.
\end{proof}

\begin{theorem}[The electron g-2]
\label{c11:thm:electron-g2}
The electron g-2 is derived from the spectral structure:
\begin{equation}
\label{c11:eq:electron-g2}
a_e = \frac{\alpha}{2\pi}\Bigl(1 - \Bigl(\frac{\mathrm{occ}_0}{\Sigma m}\Bigr)^2\Bigr)
= 1.159655\times10^{-3},
\qquad
\Delta a_e = \Bigl(\frac{\alpha}{2\pi}\Bigr)^3 \mathrm{span}^{1/4}
\Bigl(\frac{\mathrm{occ}_0}{\Sigma m}\Bigr)^3 = 1.86\times10^{-13},
\end{equation}
with the correction negative and the anomaly below $10^{-12}$.
\end{theorem}

\begin{proof}
The electron-g-2 origin \cite{c11:QG178} establishes the construction: the electron anomaly is
the Schwinger term modulated by the lightest-octave fraction, with the negative correction
$(1 - (\mathrm{occ}_0/\Sigma m)^2)$ and the anomaly below $10^{-12}$; the same D96
mechanism produces the muon g-2 \cite{c11:QG171}. The electron anomaly matches experiment
within $0.0003\%$.
\end{proof}

\begin{theorem}[The Majorana character and \texorpdfstring{$0\nu\beta\beta$}{0nubetabeta}]
\label{c11:thm:majorana}
The neutrino is Majorana: it is self-conjugate in its $T_3$-only channel, has no conserved
charge, and has a real mass matrix. The effective Majorana mass is
\begin{equation}
\label{c11:eq:mbb}
m_{\beta\beta} = \bigl|m_1c_{12}^2c_{13}^2 + m_2s_{12}^2c_{13}^2e^{i\alpha_2}
+ m_3s_{13}^2e^{-2i\delta_\nu}\bigr| = 2.02\times10^{-3}\ \mathrm{eV},
\end{equation}
within experimental limits and in reach of next-generation experiments.
\end{theorem}

\begin{proof}
The Majorana-origin audit \cite{c11:QG179} establishes the neutrino character: the $T_3$-only
channel is self-conjugate ($48/95$ modes), the unique $Q=0$ sector provides no conserved
charge, and the reflection automorphism gives a real mass matrix. From the D96 PMNS angles
and masses, the effective Majorana mass is $m_{\beta\beta} = 2.02\times10^{-3}$ eV, within
limits and in reach of next-generation experiments.
\end{proof}

The precision predictions --- $S,T,U$, the electron g-2, and the Majorana character ---
are therefore all derived from the D96 spectral structure; their experimental validation
remains open --- the current frontier, not a derivation gap (VALID001).

\begin{remark}[Hosted-SM caveat]
\label{c11:remark:hosted-sm}
Consistent with MONO007 \cite{c11:MONO007}, the standard-model \emph{Lagrangian} and its
dynamics are hosted --- a boundary, not a derivation claim of this chapter. Every
\emph{observable} presented here is derived from the D96 moments; the Lagrangian form
$L = -\frac{1}{4}F^a_{\mu\nu}F^{a\mu\nu} + i\bar{\psi}\gamma^\mu D_\mu\psi - m\bar{\psi}\psi$
is the actualization-flow action established in Chapter~\ref{c09:chap:qm}
(Lemma~\ref{c09:lem:qm-dynamics}), but the full SM dynamical content remains hosted. No stronger
claim is made.
\end{remark}

\section{Consequences}
\label{c11:sec:sm-consequences}

\begin{theorem}[The standard model is consistent with the canonical architecture]
\label{c11:thm:sm-consistent}
The standard model is consistent with the canonical architecture: every derived SM
observable is a read of the D96 spectral moments, with the hosted dynamics the single
referenced boundary. Dimensional observables are reported relative to calibration anchors,
not from the dimensionless D96 data alone.
\end{theorem}

\begin{proof}
The fermion sector (Theorem~\ref{c11:thm:fermion-access}), the gauge structure
(Theorem~\ref{c11:thm:gauge-structure}), the couplings (Theorem~\ref{c11:thm:couplings}), the flavor
mixing (Theorem~\ref{c11:thm:ckm}, Theorem~\ref{c11:thm:pmns}), the Higgs and weak scale
(Theorem~\ref{c11:thm:weak-higgs}), and the precision predictions
(Theorem~\ref{c11:thm:oblique}, Theorem~\ref{c11:thm:electron-g2}, Theorem~\ref{c11:thm:majorana})
are all derived from the D96 moments. The hosted SM Lagrangian is the single referenced
boundary (Remark~\ref{c11:remark:hosted-sm}).
\end{proof}

\section{Summary}
\label{c11:sec:sm-summary}

This chapter has reconstructed the standard-model observables: the fermion sectors are
spectral access reads (Theorem~\ref{c11:thm:fermion-access}); the $1+3+8$ gauge sector is
the D96 degree-12 observable sector (Theorem~\ref{c11:thm:gauge-structure}); the couplings
are spectral ratios (Theorem~\ref{c11:thm:couplings}); the CKM and PMNS matrices are
derived spectral reads (Theorem~\ref{c11:thm:ckm}, Theorem~\ref{c11:thm:pmns}); and the
weak scale, the Higgs mass, and the neutrino masses are spectral constructions
(Theorem~\ref{c11:thm:weak-higgs}, Lemma~\ref{c11:lem:neutrino-masses}). The precision
predictions --- the oblique parameters $S,T,U$, the electron g-2, and the Majorana
character --- are likewise reconstructed as observable outputs
(Theorem~\ref{c11:thm:oblique}, Theorem~\ref{c11:thm:electron-g2},
Theorem~\ref{c11:thm:majorana}), so the standard model is consistent with the canonical
architecture, with the hosted dynamics the single referenced boundary
(Theorem~\ref{c11:thm:sm-consistent}, Remark~\ref{c11:remark:hosted-sm}). Every SM
observable is therefore a spectral construction, twice-derived from the spectrum
(Chapter~\ref{c06:chap:d96}, Corollary~\ref{c06:cor:constants-derived}); the following
chapter derives cosmology from the same foundation.

\begin{thebibliography}{99}
\bibitem{c11:QG149} AT-QG Phase 149, \emph{Sector Exponents Origin}, Docs/Research.
\bibitem{c11:QG150} AT-QG Phase 150, \emph{Mode Access Origin}, Docs/Research.
\bibitem{c11:QG157} AT-QG Phase 157, \emph{Effective Access Counts}, Docs/Research.
\bibitem{c11:QG161} AT-QG Phase 161, \emph{Gauge Sector Origin}, Docs/Research.
\bibitem{c11:QG162} AT-QG Phase 162, \emph{Gauge Coupling Origin}, Docs/Research.
\bibitem{c11:QG165} AT-QG Phase 165, \emph{CKM Origin}, Docs/Research.
\bibitem{c11:QG166} AT-QG Phase 166, \emph{CKM CP Origin}, Docs/Research.
\bibitem{c11:QG167} AT-QG Phase 167, \emph{PMNS Origin}, Docs/Research.
\bibitem{c11:QG168} AT-QG Phase 168, \emph{Weak Boson Mass Origin}, Docs/Research.
\bibitem{c11:QG169} AT-QG Phase 169, \emph{Higgs Mass Origin}, Docs/Research.
\bibitem{c11:QG171} AT-QG Phase 171, \emph{Muon g-2 Origin}, Docs/Research.
\bibitem{c11:QG172} AT-QG Phase 172, \emph{Neutrino Mass Law}, Docs/Research.
\bibitem{c11:QG176} AT-QG Phase 176, \emph{Higgs Blind Reconstruction}, Docs/Research.
\bibitem{c11:QG178} AT-QG Phase 178, \emph{Electron g-2 Origin}, Docs/Research.
\bibitem{c11:QG179} AT-QG Phase 179, \emph{Majorana Origin}, Docs/Research.
\bibitem{c11:QG180} AT-QG Phase 180, \emph{Oblique Parameters Origin}, Docs/Research.
\bibitem{c11:MONO007} AT-MONO007, \emph{Referee-Readiness Resolution}, Docs/Zenodo/Monograph.
\end{thebibliography}
 after the preamble
%         that defines the theorem environments (or include via the master file).
% ======================================================================

% ---- Theorem environments (define in the master preamble if not present) ----
% \newtheorem{definition}{Definition}[chapter]
% \newtheorem{lemma}[definition]{Lemma}
% \newtheorem{theorem}[definition]{Theorem}
% \newtheorem{corollary}[definition]{Corollary}
% \newtheorem{remark}[definition]{Remark}

\chapter{The Standard Model}
\label{c11:chap:sm}

\section{Introduction}
\label{c11:sec:sm-introduction}

The preceding chapters established the foundation, the spectrum, the operator basis, and
the quantum and gravitational structure of The Actualization Theory. This chapter
reconstructs the \emph{standard-model observables} from the same foundation --- the fermion
sectors, the gauge structure, the couplings, the flavor mixing, the weak scale and the
Higgs, and a set of precision predictions --- as derived reads of the D96 spectral moments
\cite{c11:QG149,c11:QG150,c11:QG157}.

The central results are: the fermion sectors are reads of the spectral access counts
\cite{c11:QG149,c11:QG150}; the gauge structure $1+3+8$ is the degree-12 structure of the D96
sector \cite{c11:QG161}; the couplings are ratios of spectral constants \cite{c11:QG162}; the
flavor mixing matrices are octave-transition reads \cite{c11:QG165,c11:QG167}; the weak scale and
the Higgs mass are spectral constructions \cite{c11:QG168,c11:QG169}; the neutrino masses are
closed-form D96 laws \cite{c11:QG172}; and a set of precision predictions --- the oblique
parameters $S,T,U$, the electron g-2, and the Majorana character --- follow
\cite{c11:QG176,c11:QG178,c11:QG179,c11:QG180}.

Throughout, the hosted-standard-model caveat is retained: every \emph{observable} is
derived from the D96 moments, while the standard-model \emph{Lagrangian} and its dynamics
are hosted --- a boundary, not a derivation gap, per the referee-safe architecture of
MONO007 \cite{c11:MONO007}.

The dimensionless D96 laws are then calibrated to measured anchors: the weak scale $v$
for the weak/Higgs sector, the electron mass $m_e$ for the absolute quark masses, and the
standard unit conversion when a result is reported in SI or GeV units. That calibration
step fixes the physical units; it is not a new dynamical input.

\subsection{Derived Observables versus Hosted Dynamics}
\label{c11:sec:derived-vs-hosted}

\begin{quote}
\textbf{Derived:} primitives $\rightarrow$ actualization $\rightarrow$ spectrum $\rightarrow$
observables.\\
\textbf{Hosted:} SM Lagrangian, interaction dynamics, and external unit conventions.
\end{quote}

What is reconstructed here is the observable content of the standard model: masses,
couplings, mixings, and precision numbers. What is not reconstructed is the full dynamical
Lagrangian/vertex structure; that remains hosted and is treated as a boundary condition of
the monograph, not as a derived output.

\section{Fermion Sector}
\label{c11:sec:fermion-sector}

\begin{theorem}[The fermion sectors are spectral access reads]
\label{c11:thm:fermion-access}
The \emph{fermion sector} of the theory --- the set of fermion families and hierarchies ---
is a spectral access read: each sector couples to the spectrum through a distinct moment of
the multiplicity distribution, the neutral, full, doublet-occupancy, and octave-occupation
accesses being the half, first, second, and octave moments.
\end{theorem}

\begin{proof}
The mode-access origin \cite{c11:QG150} and the effective-access program \cite{c11:QG157}
establish that each sector couples through a distinct moment of the multiplicity
distribution: the neutral access is the half-moment $\Sigma\sqrt{m} = 64.08$, the full
access the first moment $\Sigma m = 95$, the doublet-occupancy access the second
moment $\Sigma m^2 = 229$, and the octave-occupation access
$\mathrm{occMom} = 1900.25$. The sector exponents follow from these access counts
\cite{c11:QG149}, and the mass ratios are functions of the exponents; since the counts are
themselves derived from the spectrum (Chapter~\ref{c06:chap:d96},
Corollary~\ref{c06:cor:moment-ladder}, Corollary~\ref{c06:cor:constants-derived}), the
fermion sectors and their mass hierarchy are spectral access reads.
\end{proof}

\section{Gauge Structure}
\label{c11:sec:gauge}

\begin{theorem}[The gauge structure is the D96 degree-12 observable sector]
\label{c11:thm:gauge-structure}
The gauge structure of the standard model, $U(1)\times SU(2)\times SU(3)$ --- the
$1+3+8$ gauge sector --- is reconstructed from the D96 sector structure: the degree-12 content of
the $C_{96}$ generator ring.
\end{theorem}

\begin{proof}
The gauge-origin audit \cite{c11:QG161} establishes that the $1+3+8$ gauge sector is the
degree-12 structure of the $C_{96}$ sector: the gauge generators are the D96 spectral
content, giving one $U(1)$, three $SU(2)$, and eight $SU(3)$ generators --- reconstructed as
an observable sector from the D96 spectrum, not imported, and introducing no new primitive.
Since the D96 spectrum is itself the derived output of the attractor
(Chapter~\ref{c06:chap:d96}, Theorem~\ref{c06:thm:d96-derived}), the gauge structure is
twice-emergent as an observable sector; the interaction dynamics remain hosted.
\end{proof}

\section{Coupling Constants}
\label{c11:sec:couplings}

\begin{theorem}[The couplings are spectral ratios]
\label{c11:thm:couplings}
The gauge couplings are derived as ratios of D96 spectral constants:
\begin{equation}
\label{c11:eq:couplings}
\frac{1}{\alpha_{\mathrm{em}}} = \Sigma m + \#d = 137,
\qquad
\alpha_{\mathrm{weak}} = \frac{3}{\Sigma m},
\qquad
\alpha_{\mathrm{strong}} = \frac{8}{\Sigma\sqrt{m}},
\qquad
\sin^2\theta_W = 0.2316.
\end{equation}
\end{theorem}

\begin{proof}
The coupling-origin audit \cite{c11:QG162} establishes the constructions: the fine-structure
inverse is the total mode count plus the doublet count, $1/\alpha_{\mathrm{em}} = 137$; the
weak and strong couplings are ratios of the spectral constants; the Weinberg angle is a
derived ratio. Each coupling is therefore a ratio of spectral constants ($\Sigma m$,
$\Sigma\sqrt{m}$, $\#d$) --- derived, not fitted: the constants are themselves derived
from the spectrum (Chapter~\ref{c06:chap:d96}, Corollary~\ref{c06:cor:constants-derived}),
and no fitted parameters or imported values enter.
\end{proof}

\section{Flavor Mixing}
\label{c11:sec:mixing}

\begin{theorem}[The CKM matrix is an octave-transition read]
\label{c11:thm:ckm}
The CKM matrix is derived as the quark mixing read of the spectral structure: the elements
are ratios of doublet, octave, and occupancy quantities.
\end{theorem}

\begin{proof}
The CKM-origin audit \cite{c11:QG165} establishes the constructions:
$V_{us} = \#d/(2\Sigma m)$, $V_{cb} = (\omega_0/\omega_2)^{\delta_d}$,
$V_{ub} = 2V_{cb}\cdot(\mathrm{occ}_0/\mathrm{occ}_2)$, with mean deviation $0.58\%$; the
CKM CP phase and the Jarlskog invariant follow from the chiral circulation \cite{c11:QG166}.
\end{proof}

\begin{theorem}[The PMNS matrix is a T3-only read]
\label{c11:thm:pmns}
The PMNS matrix is derived as the neutrino mixing read: the angles
$\theta_{12} = 33.35^\circ$, $\theta_{23} = 49.72^\circ$, $\theta_{13} = 8.34^\circ$, and
$\delta_\nu = 66.4^\circ$ are the T3-only access reads, with mean deviation $1.5\%$.
\end{theorem}

\begin{proof}
The PMNS-origin audit \cite{c11:QG167} establishes the angles and the CP phase as the
T3-only access reads of the D96 content; hence the PMNS matrix is a derived read, not an
imported input.
\end{proof}

Both the quark (CKM) and the lepton (PMNS) mixing matrices are therefore derived spectral
reads, with no imported mixing input.

\section{Higgs and Weak Scale}
\label{c11:sec:higgs}

\begin{theorem}[The weak scale and the Higgs mass]
\label{c11:thm:weak-higgs}
The weak scale and the Higgs mass are reconstructed from the spectral structure:
\begin{equation}
\label{c11:eq:higgs}
v = (\Sigma m+\#d)\ln(\mathrm{span}) = 254\ \mathrm{GeV},
\qquad
M_H = \sigma_{\mathrm{occ}}\cdot\frac{\mathrm{span}}{2} = 125.25\ \mathrm{GeV},
\end{equation}
with the $W$ and $Z$ masses following from the Weinberg projection.
\end{theorem}

\begin{proof}
The weak-boson-mass origin \cite{c11:QG168} establishes the weak scale
$v = (\Sigma m+\#d)\ln(\mathrm{span}) = 254$ GeV and the $W$ and $Z$ masses
($M_W = 80.1$, $M_Z = 91.4$ GeV) with $\rho = 1$ exactly; the Higgs-mass origin
\cite{c11:QG169} gives $M_H = \sigma_{\mathrm{occ}}\cdot(\mathrm{span}/2) = 125.25$ GeV,
matching observation within $0.003\%$; and the blind reconstruction \cite{c11:QG176}
confirms this mass from the pre-Higgs D96 content only, with no Higgs observable as input.
The D96 factors are dimensionless; the GeV unit is carried by the calibrated electroweak
scale $v$.
\end{proof}

The Higgs is therefore a collective scalar observable of the theory: its mass is a spectral
construction $M_H = \sigma_{\mathrm{occ}}\cdot(\mathrm{span}/2)$, whose constants are
derived from the spectrum (Chapter~\ref{c06:chap:d96},
Corollary~\ref{c06:cor:constants-derived}); it introduces no new primitive \cite{c11:QG169}.
Read strictly, the D96 factor is dimensionless and the GeV unit is inherited from the
calibrated electroweak mass scale. The Higgs sector observable is reconstructed; the Yukawa
and full interaction dynamics remain hosted.

\begin{lemma}[The neutrino masses are closed-form D96 laws]
\label{c11:lem:neutrino-masses}
The neutrino mass splittings are derived as closed-form D96 laws:
$\Delta m^2_{21} = (1/\Sigma\sqrt{m})^2/(\mathrm{span}/2)$ and
$\Delta m^2_{31} = \sin^2\theta_W/\Sigma m$.
\end{lemma}

\begin{proof}
The neutrino-mass origin \cite{c11:QG172} establishes the closed-form laws: the solar and
atmospheric splittings are ratios of the D96 spectral constants, giving
$\Delta m^2_{21} = 7.607\times10^{-5}$ eV$^2$ and
$\Delta m^2_{31} = 2.44\times10^{-3}$ eV$^2$, with $m_2 = 8.72$ meV,
$m_3 = 49.4$ meV, and $\Sigma m_\nu = 0.058$ eV. Hence the neutrino masses are closed-form
D96 laws, not oscillation fits.
\end{proof}

\section{Precision Predictions}
\label{c11:sec:precision}

\begin{theorem}[The oblique parameters]
\label{c11:thm:oblique}
The oblique parameters are derived from the spectral structure:
\begin{equation}
\label{c11:eq:oblique}
S = \frac{\mathrm{occ}_0}{\Sigma m} = \frac{4}{95} = 0.0421,
\qquad
T = 2S = \frac{8}{95} = 0.0842,
\qquad
U = 0,
\end{equation}
with $T = 2S$ exactly and $U = 0$ via the exact tree-level $\rho = 1$.
\end{theorem}

\begin{proof}
The oblique-origin audit \cite{c11:QG180} establishes the constructions: $S$ is the lightest-
octave fraction of the spectrum, $T = 2S$ exactly, and $U = 0$ via the exact SM tree-level
$\rho = 1$, matching the EW global fit within $5.3\%$.
\end{proof}

\begin{theorem}[The electron g-2]
\label{c11:thm:electron-g2}
The electron g-2 is derived from the spectral structure:
\begin{equation}
\label{c11:eq:electron-g2}
a_e = \frac{\alpha}{2\pi}\Bigl(1 - \Bigl(\frac{\mathrm{occ}_0}{\Sigma m}\Bigr)^2\Bigr)
= 1.159655\times10^{-3},
\qquad
\Delta a_e = \Bigl(\frac{\alpha}{2\pi}\Bigr)^3 \mathrm{span}^{1/4}
\Bigl(\frac{\mathrm{occ}_0}{\Sigma m}\Bigr)^3 = 1.86\times10^{-13},
\end{equation}
with the correction negative and the anomaly below $10^{-12}$.
\end{theorem}

\begin{proof}
The electron-g-2 origin \cite{c11:QG178} establishes the construction: the electron anomaly is
the Schwinger term modulated by the lightest-octave fraction, with the negative correction
$(1 - (\mathrm{occ}_0/\Sigma m)^2)$ and the anomaly below $10^{-12}$; the same D96
mechanism produces the muon g-2 \cite{c11:QG171}. The electron anomaly matches experiment
within $0.0003\%$.
\end{proof}

\begin{theorem}[The Majorana character and \texorpdfstring{$0\nu\beta\beta$}{0nubetabeta}]
\label{c11:thm:majorana}
The neutrino is Majorana: it is self-conjugate in its $T_3$-only channel, has no conserved
charge, and has a real mass matrix. The effective Majorana mass is
\begin{equation}
\label{c11:eq:mbb}
m_{\beta\beta} = \bigl|m_1c_{12}^2c_{13}^2 + m_2s_{12}^2c_{13}^2e^{i\alpha_2}
+ m_3s_{13}^2e^{-2i\delta_\nu}\bigr| = 2.02\times10^{-3}\ \mathrm{eV},
\end{equation}
within experimental limits and in reach of next-generation experiments.
\end{theorem}

\begin{proof}
The Majorana-origin audit \cite{c11:QG179} establishes the neutrino character: the $T_3$-only
channel is self-conjugate ($48/95$ modes), the unique $Q=0$ sector provides no conserved
charge, and the reflection automorphism gives a real mass matrix. From the D96 PMNS angles
and masses, the effective Majorana mass is $m_{\beta\beta} = 2.02\times10^{-3}$ eV, within
limits and in reach of next-generation experiments.
\end{proof}

The precision predictions --- $S,T,U$, the electron g-2, and the Majorana character ---
are therefore all derived from the D96 spectral structure; their experimental validation
remains open --- the current frontier, not a derivation gap (VALID001).

\begin{remark}[Hosted-SM caveat]
\label{c11:remark:hosted-sm}
Consistent with MONO007 \cite{c11:MONO007}, the standard-model \emph{Lagrangian} and its
dynamics are hosted --- a boundary, not a derivation claim of this chapter. Every
\emph{observable} presented here is derived from the D96 moments; the Lagrangian form
$L = -\frac{1}{4}F^a_{\mu\nu}F^{a\mu\nu} + i\bar{\psi}\gamma^\mu D_\mu\psi - m\bar{\psi}\psi$
is the actualization-flow action established in Chapter~\ref{c09:chap:qm}
(Lemma~\ref{c09:lem:qm-dynamics}), but the full SM dynamical content remains hosted. No stronger
claim is made.
\end{remark}

\section{Consequences}
\label{c11:sec:sm-consequences}

\begin{theorem}[The standard model is consistent with the canonical architecture]
\label{c11:thm:sm-consistent}
The standard model is consistent with the canonical architecture: every derived SM
observable is a read of the D96 spectral moments, with the hosted dynamics the single
referenced boundary. Dimensional observables are reported relative to calibration anchors,
not from the dimensionless D96 data alone.
\end{theorem}

\begin{proof}
The fermion sector (Theorem~\ref{c11:thm:fermion-access}), the gauge structure
(Theorem~\ref{c11:thm:gauge-structure}), the couplings (Theorem~\ref{c11:thm:couplings}), the flavor
mixing (Theorem~\ref{c11:thm:ckm}, Theorem~\ref{c11:thm:pmns}), the Higgs and weak scale
(Theorem~\ref{c11:thm:weak-higgs}), and the precision predictions
(Theorem~\ref{c11:thm:oblique}, Theorem~\ref{c11:thm:electron-g2}, Theorem~\ref{c11:thm:majorana})
are all derived from the D96 moments. The hosted SM Lagrangian is the single referenced
boundary (Remark~\ref{c11:remark:hosted-sm}).
\end{proof}

\section{Summary}
\label{c11:sec:sm-summary}

This chapter has reconstructed the standard-model observables: the fermion sectors are
spectral access reads (Theorem~\ref{c11:thm:fermion-access}); the $1+3+8$ gauge sector is
the D96 degree-12 observable sector (Theorem~\ref{c11:thm:gauge-structure}); the couplings
are spectral ratios (Theorem~\ref{c11:thm:couplings}); the CKM and PMNS matrices are
derived spectral reads (Theorem~\ref{c11:thm:ckm}, Theorem~\ref{c11:thm:pmns}); and the
weak scale, the Higgs mass, and the neutrino masses are spectral constructions
(Theorem~\ref{c11:thm:weak-higgs}, Lemma~\ref{c11:lem:neutrino-masses}). The precision
predictions --- the oblique parameters $S,T,U$, the electron g-2, and the Majorana
character --- are likewise reconstructed as observable outputs
(Theorem~\ref{c11:thm:oblique}, Theorem~\ref{c11:thm:electron-g2},
Theorem~\ref{c11:thm:majorana}), so the standard model is consistent with the canonical
architecture, with the hosted dynamics the single referenced boundary
(Theorem~\ref{c11:thm:sm-consistent}, Remark~\ref{c11:remark:hosted-sm}). Every SM
observable is therefore a spectral construction, twice-derived from the spectrum
(Chapter~\ref{c06:chap:d96}, Corollary~\ref{c06:cor:constants-derived}); the following
chapter derives cosmology from the same foundation.

\begin{thebibliography}{99}
\bibitem{c11:QG149} AT-QG Phase 149, \emph{Sector Exponents Origin}, Docs/Research.
\bibitem{c11:QG150} AT-QG Phase 150, \emph{Mode Access Origin}, Docs/Research.
\bibitem{c11:QG157} AT-QG Phase 157, \emph{Effective Access Counts}, Docs/Research.
\bibitem{c11:QG161} AT-QG Phase 161, \emph{Gauge Sector Origin}, Docs/Research.
\bibitem{c11:QG162} AT-QG Phase 162, \emph{Gauge Coupling Origin}, Docs/Research.
\bibitem{c11:QG165} AT-QG Phase 165, \emph{CKM Origin}, Docs/Research.
\bibitem{c11:QG166} AT-QG Phase 166, \emph{CKM CP Origin}, Docs/Research.
\bibitem{c11:QG167} AT-QG Phase 167, \emph{PMNS Origin}, Docs/Research.
\bibitem{c11:QG168} AT-QG Phase 168, \emph{Weak Boson Mass Origin}, Docs/Research.
\bibitem{c11:QG169} AT-QG Phase 169, \emph{Higgs Mass Origin}, Docs/Research.
\bibitem{c11:QG171} AT-QG Phase 171, \emph{Muon g-2 Origin}, Docs/Research.
\bibitem{c11:QG172} AT-QG Phase 172, \emph{Neutrino Mass Law}, Docs/Research.
\bibitem{c11:QG176} AT-QG Phase 176, \emph{Higgs Blind Reconstruction}, Docs/Research.
\bibitem{c11:QG178} AT-QG Phase 178, \emph{Electron g-2 Origin}, Docs/Research.
\bibitem{c11:QG179} AT-QG Phase 179, \emph{Majorana Origin}, Docs/Research.
\bibitem{c11:QG180} AT-QG Phase 180, \emph{Oblique Parameters Origin}, Docs/Research.
\bibitem{c11:MONO007} AT-MONO007, \emph{Referee-Readiness Resolution}, Docs/Zenodo/Monograph.
\end{thebibliography}
 after the preamble
%         that defines the theorem environments (or include via the master file).
% ======================================================================

% ---- Theorem environments (define in the master preamble if not present) ----
% \newtheorem{definition}{Definition}[chapter]
% \newtheorem{lemma}[definition]{Lemma}
% \newtheorem{theorem}[definition]{Theorem}
% \newtheorem{corollary}[definition]{Corollary}
% \newtheorem{remark}[definition]{Remark}

\chapter{The Standard Model}
\label{c11:chap:sm}

\section{Introduction}
\label{c11:sec:sm-introduction}

The preceding chapters established the foundation, the spectrum, the operator basis, and
the quantum and gravitational structure of The Actualization Theory. This chapter
reconstructs the \emph{standard-model observables} from the same foundation --- the fermion
sectors, the gauge structure, the couplings, the flavor mixing, the weak scale and the
Higgs, and a set of precision predictions --- as derived reads of the D96 spectral moments
\cite{c11:QG149,c11:QG150,c11:QG157}.

The central results are: the fermion sectors are reads of the spectral access counts
\cite{c11:QG149,c11:QG150}; the gauge structure $1+3+8$ is the degree-12 structure of the D96
sector \cite{c11:QG161}; the couplings are ratios of spectral constants \cite{c11:QG162}; the
flavor mixing matrices are octave-transition reads \cite{c11:QG165,c11:QG167}; the weak scale and
the Higgs mass are spectral constructions \cite{c11:QG168,c11:QG169}; the neutrino masses are
closed-form D96 laws \cite{c11:QG172}; and a set of precision predictions --- the oblique
parameters $S,T,U$, the electron g-2, and the Majorana character --- follow
\cite{c11:QG176,c11:QG178,c11:QG179,c11:QG180}.

Throughout, the hosted-standard-model caveat is retained: every \emph{observable} is
derived from the D96 moments, while the standard-model \emph{Lagrangian} and its dynamics
are hosted --- a boundary, not a derivation gap, per the referee-safe architecture of
MONO007 \cite{c11:MONO007}.

The dimensionless D96 laws are then calibrated to measured anchors: the weak scale $v$
for the weak/Higgs sector, the electron mass $m_e$ for the absolute quark masses, and the
standard unit conversion when a result is reported in SI or GeV units. That calibration
step fixes the physical units; it is not a new dynamical input.

\subsection{Derived Observables versus Hosted Dynamics}
\label{c11:sec:derived-vs-hosted}

\begin{quote}
\textbf{Derived:} primitives $\rightarrow$ actualization $\rightarrow$ spectrum $\rightarrow$
observables.\\
\textbf{Hosted:} SM Lagrangian, interaction dynamics, and external unit conventions.
\end{quote}

What is reconstructed here is the observable content of the standard model: masses,
couplings, mixings, and precision numbers. What is not reconstructed is the full dynamical
Lagrangian/vertex structure; that remains hosted and is treated as a boundary condition of
the monograph, not as a derived output.

\section{Fermion Sector}
\label{c11:sec:fermion-sector}

\begin{theorem}[The fermion sectors are spectral access reads]
\label{c11:thm:fermion-access}
The \emph{fermion sector} of the theory --- the set of fermion families and hierarchies ---
is a spectral access read: each sector couples to the spectrum through a distinct moment of
the multiplicity distribution, the neutral, full, doublet-occupancy, and octave-occupation
accesses being the half, first, second, and octave moments.
\end{theorem}

\begin{proof}
The mode-access origin \cite{c11:QG150} and the effective-access program \cite{c11:QG157}
establish that each sector couples through a distinct moment of the multiplicity
distribution: the neutral access is the half-moment $\Sigma\sqrt{m} = 64.08$, the full
access the first moment $\Sigma m = 95$, the doublet-occupancy access the second
moment $\Sigma m^2 = 229$, and the octave-occupation access
$\mathrm{occMom} = 1900.25$. The sector exponents follow from these access counts
\cite{c11:QG149}, and the mass ratios are functions of the exponents; since the counts are
themselves derived from the spectrum (Chapter~\ref{c06:chap:d96},
Corollary~\ref{c06:cor:moment-ladder}, Corollary~\ref{c06:cor:constants-derived}), the
fermion sectors and their mass hierarchy are spectral access reads.
\end{proof}

\section{Gauge Structure}
\label{c11:sec:gauge}

\begin{theorem}[The gauge structure is the D96 degree-12 observable sector]
\label{c11:thm:gauge-structure}
The gauge structure of the standard model, $U(1)\times SU(2)\times SU(3)$ --- the
$1+3+8$ gauge sector --- is reconstructed from the D96 sector structure: the degree-12 content of
the $C_{96}$ generator ring.
\end{theorem}

\begin{proof}
The gauge-origin audit \cite{c11:QG161} establishes that the $1+3+8$ gauge sector is the
degree-12 structure of the $C_{96}$ sector: the gauge generators are the D96 spectral
content, giving one $U(1)$, three $SU(2)$, and eight $SU(3)$ generators --- reconstructed as
an observable sector from the D96 spectrum, not imported, and introducing no new primitive.
Since the D96 spectrum is itself the derived output of the attractor
(Chapter~\ref{c06:chap:d96}, Theorem~\ref{c06:thm:d96-derived}), the gauge structure is
twice-emergent as an observable sector; the interaction dynamics remain hosted.
\end{proof}

\section{Coupling Constants}
\label{c11:sec:couplings}

\begin{theorem}[The couplings are spectral ratios]
\label{c11:thm:couplings}
The gauge couplings are derived as ratios of D96 spectral constants:
\begin{equation}
\label{c11:eq:couplings}
\frac{1}{\alpha_{\mathrm{em}}} = \Sigma m + \#d = 137,
\qquad
\alpha_{\mathrm{weak}} = \frac{3}{\Sigma m},
\qquad
\alpha_{\mathrm{strong}} = \frac{8}{\Sigma\sqrt{m}},
\qquad
\sin^2\theta_W = 0.2316.
\end{equation}
\end{theorem}

\begin{proof}
The coupling-origin audit \cite{c11:QG162} establishes the constructions: the fine-structure
inverse is the total mode count plus the doublet count, $1/\alpha_{\mathrm{em}} = 137$; the
weak and strong couplings are ratios of the spectral constants; the Weinberg angle is a
derived ratio. Each coupling is therefore a ratio of spectral constants ($\Sigma m$,
$\Sigma\sqrt{m}$, $\#d$) --- derived, not fitted: the constants are themselves derived
from the spectrum (Chapter~\ref{c06:chap:d96}, Corollary~\ref{c06:cor:constants-derived}),
and no fitted parameters or imported values enter.
\end{proof}

\section{Flavor Mixing}
\label{c11:sec:mixing}

\begin{theorem}[The CKM matrix is an octave-transition read]
\label{c11:thm:ckm}
The CKM matrix is derived as the quark mixing read of the spectral structure: the elements
are ratios of doublet, octave, and occupancy quantities.
\end{theorem}

\begin{proof}
The CKM-origin audit \cite{c11:QG165} establishes the constructions:
$V_{us} = \#d/(2\Sigma m)$, $V_{cb} = (\omega_0/\omega_2)^{\delta_d}$,
$V_{ub} = 2V_{cb}\cdot(\mathrm{occ}_0/\mathrm{occ}_2)$, with mean deviation $0.58\%$; the
CKM CP phase and the Jarlskog invariant follow from the chiral circulation \cite{c11:QG166}.
\end{proof}

\begin{theorem}[The PMNS matrix is a T3-only read]
\label{c11:thm:pmns}
The PMNS matrix is derived as the neutrino mixing read: the angles
$\theta_{12} = 33.35^\circ$, $\theta_{23} = 49.72^\circ$, $\theta_{13} = 8.34^\circ$, and
$\delta_\nu = 66.4^\circ$ are the T3-only access reads, with mean deviation $1.5\%$.
\end{theorem}

\begin{proof}
The PMNS-origin audit \cite{c11:QG167} establishes the angles and the CP phase as the
T3-only access reads of the D96 content; hence the PMNS matrix is a derived read, not an
imported input.
\end{proof}

Both the quark (CKM) and the lepton (PMNS) mixing matrices are therefore derived spectral
reads, with no imported mixing input.

\section{Higgs and Weak Scale}
\label{c11:sec:higgs}

\begin{theorem}[The weak scale and the Higgs mass]
\label{c11:thm:weak-higgs}
The weak scale and the Higgs mass are reconstructed from the spectral structure:
\begin{equation}
\label{c11:eq:higgs}
v = (\Sigma m+\#d)\ln(\mathrm{span}) = 254\ \mathrm{GeV},
\qquad
M_H = \sigma_{\mathrm{occ}}\cdot\frac{\mathrm{span}}{2} = 125.25\ \mathrm{GeV},
\end{equation}
with the $W$ and $Z$ masses following from the Weinberg projection.
\end{theorem}

\begin{proof}
The weak-boson-mass origin \cite{c11:QG168} establishes the weak scale
$v = (\Sigma m+\#d)\ln(\mathrm{span}) = 254$ GeV and the $W$ and $Z$ masses
($M_W = 80.1$, $M_Z = 91.4$ GeV) with $\rho = 1$ exactly; the Higgs-mass origin
\cite{c11:QG169} gives $M_H = \sigma_{\mathrm{occ}}\cdot(\mathrm{span}/2) = 125.25$ GeV,
matching observation within $0.003\%$; and the blind reconstruction \cite{c11:QG176}
confirms this mass from the pre-Higgs D96 content only, with no Higgs observable as input.
The D96 factors are dimensionless; the GeV unit is carried by the calibrated electroweak
scale $v$.
\end{proof}

The Higgs is therefore a collective scalar observable of the theory: its mass is a spectral
construction $M_H = \sigma_{\mathrm{occ}}\cdot(\mathrm{span}/2)$, whose constants are
derived from the spectrum (Chapter~\ref{c06:chap:d96},
Corollary~\ref{c06:cor:constants-derived}); it introduces no new primitive \cite{c11:QG169}.
Read strictly, the D96 factor is dimensionless and the GeV unit is inherited from the
calibrated electroweak mass scale. The Higgs sector observable is reconstructed; the Yukawa
and full interaction dynamics remain hosted.

\begin{lemma}[The neutrino masses are closed-form D96 laws]
\label{c11:lem:neutrino-masses}
The neutrino mass splittings are derived as closed-form D96 laws:
$\Delta m^2_{21} = (1/\Sigma\sqrt{m})^2/(\mathrm{span}/2)$ and
$\Delta m^2_{31} = \sin^2\theta_W/\Sigma m$.
\end{lemma}

\begin{proof}
The neutrino-mass origin \cite{c11:QG172} establishes the closed-form laws: the solar and
atmospheric splittings are ratios of the D96 spectral constants, giving
$\Delta m^2_{21} = 7.607\times10^{-5}$ eV$^2$ and
$\Delta m^2_{31} = 2.44\times10^{-3}$ eV$^2$, with $m_2 = 8.72$ meV,
$m_3 = 49.4$ meV, and $\Sigma m_\nu = 0.058$ eV. Hence the neutrino masses are closed-form
D96 laws, not oscillation fits.
\end{proof}

\section{Precision Predictions}
\label{c11:sec:precision}

\begin{theorem}[The oblique parameters]
\label{c11:thm:oblique}
The oblique parameters are derived from the spectral structure:
\begin{equation}
\label{c11:eq:oblique}
S = \frac{\mathrm{occ}_0}{\Sigma m} = \frac{4}{95} = 0.0421,
\qquad
T = 2S = \frac{8}{95} = 0.0842,
\qquad
U = 0,
\end{equation}
with $T = 2S$ exactly and $U = 0$ via the exact tree-level $\rho = 1$.
\end{theorem}

\begin{proof}
The oblique-origin audit \cite{c11:QG180} establishes the constructions: $S$ is the lightest-
octave fraction of the spectrum, $T = 2S$ exactly, and $U = 0$ via the exact SM tree-level
$\rho = 1$, matching the EW global fit within $5.3\%$.
\end{proof}

\begin{theorem}[The electron g-2]
\label{c11:thm:electron-g2}
The electron g-2 is derived from the spectral structure:
\begin{equation}
\label{c11:eq:electron-g2}
a_e = \frac{\alpha}{2\pi}\Bigl(1 - \Bigl(\frac{\mathrm{occ}_0}{\Sigma m}\Bigr)^2\Bigr)
= 1.159655\times10^{-3},
\qquad
\Delta a_e = \Bigl(\frac{\alpha}{2\pi}\Bigr)^3 \mathrm{span}^{1/4}
\Bigl(\frac{\mathrm{occ}_0}{\Sigma m}\Bigr)^3 = 1.86\times10^{-13},
\end{equation}
with the correction negative and the anomaly below $10^{-12}$.
\end{theorem}

\begin{proof}
The electron-g-2 origin \cite{c11:QG178} establishes the construction: the electron anomaly is
the Schwinger term modulated by the lightest-octave fraction, with the negative correction
$(1 - (\mathrm{occ}_0/\Sigma m)^2)$ and the anomaly below $10^{-12}$; the same D96
mechanism produces the muon g-2 \cite{c11:QG171}. The electron anomaly matches experiment
within $0.0003\%$.
\end{proof}

\begin{theorem}[The Majorana character and \texorpdfstring{$0\nu\beta\beta$}{0nubetabeta}]
\label{c11:thm:majorana}
The neutrino is Majorana: it is self-conjugate in its $T_3$-only channel, has no conserved
charge, and has a real mass matrix. The effective Majorana mass is
\begin{equation}
\label{c11:eq:mbb}
m_{\beta\beta} = \bigl|m_1c_{12}^2c_{13}^2 + m_2s_{12}^2c_{13}^2e^{i\alpha_2}
+ m_3s_{13}^2e^{-2i\delta_\nu}\bigr| = 2.02\times10^{-3}\ \mathrm{eV},
\end{equation}
within experimental limits and in reach of next-generation experiments.
\end{theorem}

\begin{proof}
The Majorana-origin audit \cite{c11:QG179} establishes the neutrino character: the $T_3$-only
channel is self-conjugate ($48/95$ modes), the unique $Q=0$ sector provides no conserved
charge, and the reflection automorphism gives a real mass matrix. From the D96 PMNS angles
and masses, the effective Majorana mass is $m_{\beta\beta} = 2.02\times10^{-3}$ eV, within
limits and in reach of next-generation experiments.
\end{proof}

The precision predictions --- $S,T,U$, the electron g-2, and the Majorana character ---
are therefore all derived from the D96 spectral structure; their experimental validation
remains open --- the current frontier, not a derivation gap (VALID001).

\begin{remark}[Hosted-SM caveat]
\label{c11:remark:hosted-sm}
Consistent with MONO007 \cite{c11:MONO007}, the standard-model \emph{Lagrangian} and its
dynamics are hosted --- a boundary, not a derivation claim of this chapter. Every
\emph{observable} presented here is derived from the D96 moments; the Lagrangian form
$L = -\frac{1}{4}F^a_{\mu\nu}F^{a\mu\nu} + i\bar{\psi}\gamma^\mu D_\mu\psi - m\bar{\psi}\psi$
is the actualization-flow action established in Chapter~\ref{c09:chap:qm}
(Lemma~\ref{c09:lem:qm-dynamics}), but the full SM dynamical content remains hosted. No stronger
claim is made.
\end{remark}

\section{Consequences}
\label{c11:sec:sm-consequences}

\begin{theorem}[The standard model is consistent with the canonical architecture]
\label{c11:thm:sm-consistent}
The standard model is consistent with the canonical architecture: every derived SM
observable is a read of the D96 spectral moments, with the hosted dynamics the single
referenced boundary. Dimensional observables are reported relative to calibration anchors,
not from the dimensionless D96 data alone.
\end{theorem}

\begin{proof}
The fermion sector (Theorem~\ref{c11:thm:fermion-access}), the gauge structure
(Theorem~\ref{c11:thm:gauge-structure}), the couplings (Theorem~\ref{c11:thm:couplings}), the flavor
mixing (Theorem~\ref{c11:thm:ckm}, Theorem~\ref{c11:thm:pmns}), the Higgs and weak scale
(Theorem~\ref{c11:thm:weak-higgs}), and the precision predictions
(Theorem~\ref{c11:thm:oblique}, Theorem~\ref{c11:thm:electron-g2}, Theorem~\ref{c11:thm:majorana})
are all derived from the D96 moments. The hosted SM Lagrangian is the single referenced
boundary (Remark~\ref{c11:remark:hosted-sm}).
\end{proof}

\section{Summary}
\label{c11:sec:sm-summary}

This chapter has reconstructed the standard-model observables: the fermion sectors are
spectral access reads (Theorem~\ref{c11:thm:fermion-access}); the $1+3+8$ gauge sector is
the D96 degree-12 observable sector (Theorem~\ref{c11:thm:gauge-structure}); the couplings
are spectral ratios (Theorem~\ref{c11:thm:couplings}); the CKM and PMNS matrices are
derived spectral reads (Theorem~\ref{c11:thm:ckm}, Theorem~\ref{c11:thm:pmns}); and the
weak scale, the Higgs mass, and the neutrino masses are spectral constructions
(Theorem~\ref{c11:thm:weak-higgs}, Lemma~\ref{c11:lem:neutrino-masses}). The precision
predictions --- the oblique parameters $S,T,U$, the electron g-2, and the Majorana
character --- are likewise reconstructed as observable outputs
(Theorem~\ref{c11:thm:oblique}, Theorem~\ref{c11:thm:electron-g2},
Theorem~\ref{c11:thm:majorana}), so the standard model is consistent with the canonical
architecture, with the hosted dynamics the single referenced boundary
(Theorem~\ref{c11:thm:sm-consistent}, Remark~\ref{c11:remark:hosted-sm}). Every SM
observable is therefore a spectral construction, twice-derived from the spectrum
(Chapter~\ref{c06:chap:d96}, Corollary~\ref{c06:cor:constants-derived}); the following
chapter derives cosmology from the same foundation.

\begin{thebibliography}{99}
\bibitem{c11:QG149} AT-QG Phase 149, \emph{Sector Exponents Origin}, Docs/Research.
\bibitem{c11:QG150} AT-QG Phase 150, \emph{Mode Access Origin}, Docs/Research.
\bibitem{c11:QG157} AT-QG Phase 157, \emph{Effective Access Counts}, Docs/Research.
\bibitem{c11:QG161} AT-QG Phase 161, \emph{Gauge Sector Origin}, Docs/Research.
\bibitem{c11:QG162} AT-QG Phase 162, \emph{Gauge Coupling Origin}, Docs/Research.
\bibitem{c11:QG165} AT-QG Phase 165, \emph{CKM Origin}, Docs/Research.
\bibitem{c11:QG166} AT-QG Phase 166, \emph{CKM CP Origin}, Docs/Research.
\bibitem{c11:QG167} AT-QG Phase 167, \emph{PMNS Origin}, Docs/Research.
\bibitem{c11:QG168} AT-QG Phase 168, \emph{Weak Boson Mass Origin}, Docs/Research.
\bibitem{c11:QG169} AT-QG Phase 169, \emph{Higgs Mass Origin}, Docs/Research.
\bibitem{c11:QG171} AT-QG Phase 171, \emph{Muon g-2 Origin}, Docs/Research.
\bibitem{c11:QG172} AT-QG Phase 172, \emph{Neutrino Mass Law}, Docs/Research.
\bibitem{c11:QG176} AT-QG Phase 176, \emph{Higgs Blind Reconstruction}, Docs/Research.
\bibitem{c11:QG178} AT-QG Phase 178, \emph{Electron g-2 Origin}, Docs/Research.
\bibitem{c11:QG179} AT-QG Phase 179, \emph{Majorana Origin}, Docs/Research.
\bibitem{c11:QG180} AT-QG Phase 180, \emph{Oblique Parameters Origin}, Docs/Research.
\bibitem{c11:MONO007} AT-MONO007, \emph{Referee-Readiness Resolution}, Docs/Zenodo/Monograph.
\end{thebibliography}
